\documentclass[conference]{IEEEtran}
\IEEEoverridecommandlockouts

\usepackage{cite}
\usepackage{amsmath,amssymb}
\usepackage{graphicx}
\usepackage{booktabs}
\usepackage{tabularx}
\usepackage{array}
\usepackage{xcolor}
\usepackage{tikz}
\usetikzlibrary{positioning,arrows.meta,decorations.pathreplacing,calc}
\usepackage[hidelinks]{hyperref}

\newcommand{\todo}[1]{\textcolor{red}{[TODO: #1]}}
\newcommand{\tsens}{t_{\text{sensor}}}
\newcommand{\tos}{t_{\text{os}}}
\newcommand{\tqueue}{t_{\text{queue}}}
\newcommand{\tdone}{t_{\text{done}}}

\begin{document}

\title{Autonomy Is Not Free Attention}

\author{\IEEEauthorblockN{Yojan Gautam}
\IEEEauthorblockA{George Mason University\\
ygautam2@gmu.edu}
}

\maketitle

\begin{abstract}


\end{abstract}

\begin{IEEEkeywords}
human–robot interaction, supervisory control, multi-robot systems, operator workload, attention allocation
\end{IEEEkeywords}

% =====================================================================
\section{Introduction}
\label{sec:intro}
The desire to delegate work to machines predates computing itself.
Wiener's cybernetics framed machines as systems that could close
their own control loops through feedback, freeing the human from
moment-to-moment regulation~\cite{wiener1948cybernetics}. As machines
became capable of acting on their own, the human's role shifted from
\emph{operator} to \emph{supervisor}: someone who sets goals,
monitors progress, and intervenes when the machine cannot cope
\cite{sheridan1992telerobotics}. Autonomy is now commonly described
not as a binary property but as a spectrum of levels and stages,
each allocating a different share of sensing, decision, and action
to the machine~\cite{parasuraman2000model}.

This shift carries an appealing promise. If each agent needs less
of its operator, then one operator should be able to supervise more
agents. Olsen and Wood formalized this idea as \emph{fan-out}, the
number of robots a single human can effectively control, and showed
that it grows with how long a robot can be left unattended relative
to how long it takes to service~\cite{olsen2004fanout}. Crandall
\emph{et al.}\ grounded this in the notions of \emph{neglect time}
(NT) and \emph{interaction time} (IT)~\cite{crandall2005validating},
and later work extended the model to account for the time robots
spend waiting for an operator who is busy
elsewhere~\cite{cummings2008predicting} and to identify which metrics
predict team performance~\cite{crandall2007identifying}. Recent field
studies suggest that, with sufficiently capable autonomy, a single
operator can deploy a swarm of 100 heterogeneous ground and aerial
robots~\cite{adams2023single}, and new planning
approaches treat the operator explicitly as a scarce resource to be
allocated across a fleet~\cite{costen2024multirobot}. The implicit
assumption across much of this work is that more autonomy yields
longer neglect time and therefore higher fan-out: attention becomes
cheaper as agents become more capable.

Human factors research has long cautioned against this assumption.
Bainbridge's \emph{ironies of automation} observed that automating a
task does not remove the human; it leaves them with the hardest
residual work, monitoring a system they no longer actively control
and taking over precisely when it fails~\cite{bainbridge1983ironies}.
Operators removed from the control loop lose situation awareness and
respond more slowly and less accurately when they must
re-engage~\cite{endsley1995outoftheloop}. A meta-analysis across
studies that varied the degree of automation found that higher
automation improves routine performance and reduces workload, but
degrades situation awareness and performance when the automation
fails~\cite{onnasch2014human}, a pattern summarized as the
``lumberjack effect.'' Partially automated driving offers a familiar
illustration: at SAE Level~2, the vehicle steers and controls speed,
yet the driver remains responsible for continuously supervising it
and taking over immediately when needed~\cite{sae2021j3016}. The
work of driving is not eliminated; it becomes the work of watching.

A key reason this residual work is costly is that higher autonomy
hides more of the agent's decision-making. When a capable agent acts
on its own, the operator often cannot see \emph{why} it did what it
did. Intervening effectively then requires first reconstructing the
agent's current plan, its reasoning, and its expected outcome, the
three levels of information that the Situation Awareness-based Agent
Transparency (SAT) model identifies as necessary for effective
human--agent teaming~\cite{chen2018situation}. How much that
reconstruction costs depends on who the operator is. An engineer who
built the system carries a mental model of its planner, its failure
modes, and its quirks, and can often infer what an agent is doing
from a few cues. The people expected to supervise deployed fleets,
such as first responders or field technicians, usually lack that
model and must rebuild the agent's reasoning from what the interface
shows them. Existing evidence for high fan-out does not test this
operator over a full shift. The field demonstration in which one
commander deployed a swarm of 100 heterogeneous robots relied on two
core team members who were also the system's
developers~\cite{adams2023single}. Laboratory studies, by contrast,
recruit mostly students and other
novices~\cite{hocraffer2017metaanalysis,reybecerra2026mastering,crandall2011computing},
but test them in simplified simulations for sessions of 10 to 30
minutes~\cite{clare2010assessing,cummings2008predicting}, in one case
compressing a strike that would take several hours into half an
hour~\cite{cummings2008predicting}. Neither setting measures how many
agents a person who did not build the system can supervise over a
full working shift.

Session length matters because sustained monitoring is effortful.
Vigilance declines with time on watch~\cite{warm2008vigilance}, and higher autonomy
can itself induce fatigue-like states of monotony and reduced
vigilance, because the operator is left with less to
do~\cite{reybecerra2026mastering}. Even the developers in the field
demonstration above reported fatigue rising within shifts of 1 to 3.5
hours, and the two commanders traded shifts to sustain the
exercise~\cite{adams2023single}. Operations that outlast one
operator's endurance require rotation, and each incoming operator must
rebuild an understanding of every agent in the fleet. Handoffs in
high-consequence supervisory settings demand deliberate coordination
strategies~\cite{patterson2004handoff}, and operators handing off
vehicles in multi-vehicle operations benefit from dedicated tools to
manage the transfer~\cite{chandarana2022streamlining}; yet the human
factors of handover in uncrewed systems remain sparsely
studied~\cite{grindley2025handover}. A
fan-out estimate taken from a short session therefore overstates what
a fleet can sustain in practice.

Without transparency, the operator must either rely on the agent
without verifying it or spend attention verifying it. Trust governs
which of these happens: well-calibrated trust is necessary for
appropriate reliance~\cite{lee2004trust}, whereas over-trust leads to
misuse and missed failures, and under-trust leads to disuse and
unnecessary monitoring~\cite{parasuraman1997humans}. Trust is also
dynamic; it rises and falls as operators observe agent successes and
failures over time~\cite{hoff2015trust,devisser2020towards}, so the
attention an operator devotes to an agent is not a fixed property of
its autonomy level.

These observations suggest that the relationship between autonomy
and fan-out is more complex than classic models assume. Increasing
autonomy may lengthen neglect time, but it may also lengthen
interaction time, because each intervention now begins with the
cost of understanding an opaque agent and recalibrating trust in it,
a cost that is larger for operators who did not build the system and
that grows as they tire. Trust adds two further costs. First, an
operator who distrusts a highly autonomous agent under-relies on it,
checking on it and stepping in when nothing is
wrong~\cite{parasuraman1997humans,lee2004trust}. This \emph{trust tax}
consumes operator time that the fan-out equation does not count and
cuts short the agent's effective neglect time. Second, an operator who
stops attending to an autonomous agent pays a \emph{recovery cost}
when the autonomy gets stuck: the failure waits longer to be
noticed~\cite{cummings2008predicting}, and the interaction needed to
bring the agent back above its performance threshold depends on how
far its state has degraded, which grows the longer it was
neglected~\cite{cummings2008predicting,crandall2011computing}. If
these costs grow faster than the gains in neglect time, fan-out will
\emph{not} increase monotonically with autonomy.

Classic fan-out models also treat each robot independently: the time
an operator spends on one robot does not change what any other robot
costs. Human factors research suggests otherwise. The operator's
situation awareness is a single, shared resource, and it depends on
how the operator is engaged. Operators working at intermediate levels
of automation keep better situation awareness and recover from
automation failures faster than operators who only
monitor~\cite{endsley1995outoftheloop,kaber2004effects}, and handing
a task back to the operator for a while restores their ability to
detect failures in the automation that remains~\cite{parasuraman1996effects}.
In a fleet, this means one rover's autonomy level can affect the
cost of supervising the others. An operator who is actively directing
some rovers stays in the loop. When an autonomous rover gets stuck,
that operator already knows the terrain, the mission's progress, and
what the fleet is doing, and has less to reconstruct before acting.
A fleet that keeps some rovers under closer human control may
therefore make its autonomous rovers cheaper to recover. Low-autonomy
rovers are not free either, though: teleoperation consumes nearly all
of the operator's attention while it lasts. Whether a mixed fleet
supports more rovers than a uniform one depends on whether the
recovery time saved on autonomous rovers exceeds the interaction time
spent on closely controlled ones, and that balance changes with the
task, the terrain, and the operator's state. Choosing the mix is
therefore an optimization problem, not a design constant. Despite
repeated calls to determine what level of autonomy robots in a
supervised team should have and when that level should
change~\cite{crandall2007identifying,chen2014human}, this trade-off
has received little direct empirical study in multi-robot settings.

In this paper, we argue that \emph{autonomy is not free attention}:
autonomy does not remove the attentional cost of supervision but
shifts it from \emph{controlling} agents to \emph{verifying} them,
and an operator who is out of the loop pays that cost at every
recovery. We claim that a fleet in which an optimizer assigns each
rover its level of autonomy, keeping some rovers under closer human
control, can be supervised at a higher fan-out than a fleet at any
single uniform level, even though fewer of its rovers are fully
autonomous. We attribute this gain to two sources and measure them
separately. The first is \emph{allocation}: matching each rover's
level to the difficulty of its task. The second is \emph{engagement}:
hands-on control keeps the operator's situation awareness high
enough that recovering the autonomous rovers takes less time. We
test this claim in a user study in which a single operator supervises
a fleet of ground rovers under uniform autonomy levels, a fixed mix
of levels, and a mix recommended by the optimizer.
\todo{State participant pool (non-developers?) and session length
here, so the study matches the motivation above.}
Our contributions are:

\begin{itemize}
  \item A fan-out model for fleets with \emph{mixed} levels of
        autonomy. It replaces the per-level fan-out equation with a
        constraint on operator utilization summed over rovers, and
        extends neglect, interaction, and wait time (NT/IT/WT) with
        (a) a \emph{trust tax}, the monitoring and unprompted
        interventions spent on rovers that do not need help;
        (b) a \emph{recovery cost} that grows with the time since a
        rover last received attention; and (c) a fleet-level
        \emph{operator engagement} state, raised by active control and
        decaying during passive monitoring, that scales the
        interaction time of every rover. Terms (a) and (b) are tied to
        a latent per-rover trust state with explicit dynamics.
  \item An autonomy allocation optimizer that uses this model to
        recommend a level for each rover, accounting for the effect
        each rover's level has on the cost of supervising the rest of
        the fleet.
  \item A multi-rover supervision platform supporting five levels of
        autonomy, from teleoperation to exploration (L0--L4,
        Table~\ref{tab:levels}), with per-rover logging of neglect,
        interaction, and wait time and in-task probes of situation
        awareness.
  \item A user study comparing uniform autonomy levels, a fixed mix,
        and the optimizer's recommended mix on fan-out, task
        performance, situation awareness, operator workload
        (NASA-TLX~\cite{hart1988development}), and trust. Comparing
        the fixed and recommended mixes separates the gain from
        mixing levels from the gain from optimizing them, and
        comparing the recovery time of autonomous rovers in mixed and
        uniform fleets tests the engagement mechanism directly.
        \todo{Summarize key finding once data is in.}
\end{itemize}


\section{Problem Formulation}
\label{sec:problem}

\subsection{Levels of Autonomy}
\label{sec:levels}
Each rover runs at one of five levels of autonomy, L0 to L4
(Table~\ref{tab:levels}). The levels differ in who sets the rover's
goals and who is responsible for avoiding obstacles. At L0 and L1 the
operator drives; at L2 and L3 the operator places waypoints and the
rover drives to them; at L4 the rover chooses its own goals. Each step
up hands the operator's next task to the rover: L1 takes over
stopping for obstacles, L2 takes over driving, L3 takes over avoiding
obstacles, and L4 takes over choosing where to go.

At L3 the rover avoids an obstacle only when it can find a path around
it. When it cannot, it stops and signals the operator, who must then
intervene. L3 also requires every waypoint to lie inside a map the
rover has already built, either while the operator drove it manually
or during L2 waypoint runs. Only L4 can take the rover into
locations it has not mapped.

\begin{table}[t]
\caption{Levels of autonomy. Each level hands one more task from the
operator to the rover.}
\label{tab:levels}
\centering
\small
\begin{tabularx}{\columnwidth}{@{}l>{\raggedright\arraybackslash}X>{\raggedright\arraybackslash}X@{}}
\toprule
Level & Operator & Rover \\
\midrule
L0 & Teleoperates; responsible for avoiding obstacles. & Nothing on its own. \\
L1 & Teleoperates. & Recognizes obstacles and stops. \\
L2 & Sets waypoints; monitors for obstacles. & Drives to each waypoint; does not avoid obstacles. \\
L3 & Sets waypoints inside the existing map. & Drives to each waypoint and avoids obstacles; if it finds no path, stops and signals the operator. \\
L4 & Monitors; takes over when needed. & Explores new locations, plans paths, and avoids obstacles. \\
\bottomrule
\end{tabularx}
\end{table}

\subsection{Classic Fan-Out}
\label{sec:classic-fo}
Each rover alternates between two phases. During a \emph{neglect}
interval it works without the operator; the interval ends when the
rover needs help, either because it has finished its command and needs
a new one or because it has failed. The rover then \emph{waits} until
the operator turns to it, and an \emph{interaction} follows in which
the operator gives a new command or recovers the rover. With mean
neglect time $\mathrm{NT}_\ell$, interaction time $\mathrm{IT}_\ell$,
and wait time $\mathrm{WT}$ at level $\ell$, the number of rovers one
operator can keep productive
is~\cite{olsen2004fanout,cummings2008predicting}
\begin{equation}
  \mathrm{FO}_\ell = \frac{\mathrm{NT}_\ell}{\mathrm{IT}_\ell + \mathrm{WT}} + 1 .
  \label{eq:fo-classic}
\end{equation}
Higher levels lengthen $\mathrm{NT}_\ell$, so Eq.~\eqref{eq:fo-classic}
predicts that fan-out rises with autonomy. The equation counts only
operator time spent in interactions. Time spent watching a rover that
does not need help appears nowhere, and $\mathrm{IT}_\ell$ is a
constant that does not depend on how long the rover went unattended.
We relax both assumptions through a model of operator trust
(Sections~\ref{sec:trust}--\ref{sec:fo-aware}). The equation also
assumes that every rover runs at the same level and that what one
rover costs the operator does not depend on the others. We relax
these assumptions through a model of operator engagement
(Sections~\ref{sec:engagement}--\ref{sec:allocation}).

\subsection{Operator Trust}
\label{sec:trust}
Let $\theta_i \in [0,1]$ be the operator's trust in rover $i$ at its
current level and in its current situation. We track trust separately
for each level and situation, such as terrain type, because an
operator may trust a rover at L3 on grass but not on rubble.
Trust is latent: we estimate it from questionnaires and from the
operator's behavior, but we do not observe it directly. It changes
when the operator observes an outcome. Following the linear dynamic
trust model of Yao and Zhang~\cite{yao2020optimal}, after the $k$-th
observed outcome of rover $i$,
\begin{equation}
  \theta_i(k) = \Big[\, A\,\theta_i(k-1) + B\,P_i(k) - C\,P_i(k-1) \,\Big]_0^1 ,
  \label{eq:trust}
\end{equation}
where $[\cdot]_0^1$ clips trust to $[0,1]$. The coefficient
$A \in (0,1)$ is memory: values just under one mean trust changes
slowly. The term $B\,P_i(k)$ raises trust when the rover performs well,
and $-C\,P_i(k-1)$ makes trust respond to changes in performance as
well as to its level. We fit $A$, $B$, and $C$ for each participant.
The performance score $P_i(k)$ comes from logged events:
\begin{equation}
  P_i(k) =
  \begin{cases}
    p_s  & \text{command completed without help,}\\
    -p_a & \text{announced failure,}\\
    -p_q & \text{silent failure the operator discovered,}
  \end{cases}
  \label{eq:performance}
\end{equation}
with $p_s, p_a, p_q > 0$. We expect $p_q > p_a$: a failure the operator
was not warned about should cost more trust than one the rover
announced~\cite{lee2004trust}. This is consistent with the observation
that trust rises and falls with experience of the
agent~\cite{hoff2015trust}. Only \emph{observed} outcomes update
trust. A failure the operator never notices leaves trust unchanged,
so an operator who rarely checks a rover can keep trusting it while
it is stuck.

We measure trust in two ways. Participants rate their trust in each
rover with a short scale after each trial and with an in-task slider
during it. \todo{Choose and cite the scale, e.g., Jian et al.\ or
Muir.} We fit $A$, $B$, and $C$ to these ratings rather than inferring
trust from performance alone, since reported trust need not track
performance. Behavior gives the second measure: the rate at which an
operator checks rovers that do not need help, and steps in when they
did not have to. The trust tax (Section~\ref{sec:trust-tax}) rests on
the assumption that the rated trust predicts this behavior, and we
test that assumption directly.

\subsection{Trust Tax}
\label{sec:trust-tax}
While rover $i$ works at level $\ell$, the operator checks on it at
rate $c_\ell(\theta_i)$, each check taking mean time $\delta$, and
steps in without being asked at rate $\nu_\ell(\theta_i)$, each
unprompted intervention taking mean time $\mathrm{IT}^{u}_\ell$. Both
rates fall as trust rises; in the simplest form they are linear in
distrust, $c_\ell(\theta) = c^{0}_\ell + c^{1}_\ell\,(1-\theta)$ and
$\nu_\ell(\theta) = \nu^{1}_\ell\,(1-\theta)$, where $c^{0}_\ell$ is the
check rate of a fully trusting operator. A check sends no command, so it is not
logged as interaction time. An unprompted intervention ends the
neglect interval early, so it competes with the rover's own need for
help. If the rover needs help at rate $1/\mathrm{NT}_\ell$, its
effective neglect time is
\begin{equation}
  \mathrm{NT}^{\mathrm{eff}}_i =
  \left( \frac{1}{\mathrm{NT}_\ell} + \nu_\ell(\theta_i) \right)^{-1}
  \le \mathrm{NT}_\ell .
  \label{eq:nt-eff}
\end{equation}
The \emph{trust tax} is the operator time rover $i$ takes per cycle
while it does not need help:
\begin{equation}
  \tau_i = \mathrm{NT}^{\mathrm{eff}}_i
  \left( c_\ell(\theta_i)\,\delta + \nu_\ell(\theta_i)\,\mathrm{IT}^{u}_\ell \right).
  \label{eq:trust-tax}
\end{equation}

\subsection{Recovery Cost}
\label{sec:recovery}
A fraction $h_\ell$ of neglect intervals end in failure rather than
in a completed command. Of these failures, a fraction $\varphi_\ell$
are \emph{announced}: the rover stops and signals the operator, as an
L3 rover does when it finds no path. The rest are \emph{silent}, and
the operator learns of them only at the next check. If checks arrive
at rate $c_\ell(\theta_i)$, a silent failure waits on average
$1/c_\ell(\theta_i)$ before it is noticed, so the expected wait per
cycle is
\begin{equation}
  \mathrm{WT}_i = W_q + h_\ell\,\frac{1-\varphi_\ell}{c_\ell(\theta_i)},
  \label{eq:wt}
\end{equation}
where $W_q$ is the delay while the operator serves other rovers.

Let $a_i$ be the time since rover $i$ last received the operator's
attention, whether a command, a check, or an interaction, measured
when recovery begins. The longer a stuck rover goes unattended, the
more its state can degrade and the more the operator must reconstruct
before acting. We model the interaction time needed to bring it back
above its performance threshold as
\begin{equation}
  \mathrm{IT}^{\mathrm{rec}}_\ell(a_i) = \mathrm{IT}^{\mathrm{rec}}_{\ell,0} + \beta_\ell\, a_i,
  \qquad \beta_\ell \ge 0 ,
  \label{eq:it-rec}
\end{equation}
where the \emph{recovery cost} is the term $\beta_\ell a_i$. The
expected interaction time per cycle mixes commands and recoveries:
\begin{equation}
  \mathrm{IT}^{\mathrm{eff}}_i = (1-h_\ell)\,\mathrm{IT}^{\mathrm{cmd}}_\ell
  + h_\ell\,\mathbb{E}\!\left[\mathrm{IT}^{\mathrm{rec}}_\ell(a_i)\right].
  \label{eq:it-eff}
\end{equation}

\subsection{Trust- and Recovery-Aware Fan-Out}
\label{sec:fo-aware}
Counting the trust tax as operator time that each rover consumes, the
fan-out of a fleet at level $\ell$ becomes
\begin{equation}
  \mathrm{FO}_\ell(\theta) =
  \frac{\mathrm{NT}^{\mathrm{eff}}}{\mathrm{IT}^{\mathrm{eff}} + \mathrm{WT} + \tau} + 1 .
  \label{eq:fo-aware}
\end{equation}
The trust tax lowers fan-out twice, by shortening
$\mathrm{NT}^{\mathrm{eff}}$ and by adding $\tau$ to the denominator.
The recovery cost lowers it by raising $\mathrm{IT}^{\mathrm{eff}}$ and
$\mathrm{WT}$.

Eq.~\eqref{eq:fo-aware} reduces to Eq.~\eqref{eq:fo-classic} when the
operator never checks or intervenes unprompted ($c_\ell = \nu_\ell =
0$), recovery does not depend on neglect ($\beta_\ell = 0$), and every
failure is announced ($\varphi_\ell = 1$). We call this the
\emph{naive model}. It treats autonomy as free attention, and it is
the baseline against which we compare the full model.

The check rate $c_\ell(\theta)$ enters both costs with opposite sign.
Checking more raises the trust tax; checking less delays the discovery
of silent failures and lengthens the time each stuck rover goes
unattended, which raises the recovery cost. Fan-out is therefore
highest at an intermediate check rate. The operator does not choose
this rate directly: trust sets it. An operator who under-trusts a
rover checks it more than needed and pays the trust tax; one who
over-trusts it checks it too little and pays the recovery cost. The
model gives calibrated trust a measurable meaning, namely the check
rate that maximizes Eq.~\eqref{eq:fo-aware}. It also shows how fan-out
can fall as autonomy rises. A higher level lengthens
$\mathrm{NT}_\ell$, but if it hides more of the rover's reasoning, it
may also raise the check rate at a given trust, make more failures
silent, or raise $\beta_\ell$, and any of these can outweigh the gain
in neglect time.

\subsection{Operator Engagement}
\label{sec:engagement}
The model so far gives each rover its own trust state and its own
costs. The operator's understanding of the fleet as a whole, however,
is a single shared state. Let $E(t) \in [0,1]$ be the operator's
\emph{engagement}: how well the operator currently knows the terrain,
the progress of the mission, and what the fleet is doing. Like trust,
engagement is latent; we estimate it from situation awareness probes
during the task. Engagement rises while the operator actively controls
a rover and decays while the operator only
monitors~\cite{endsley1995outoftheloop,kaber2004effects}.

Let $g_\ell \in [0,1]$ be how much active control an interaction at
level $\ell$ involves, with $g_\ell$ decreasing in $\ell$. Driving a
rover at L0 or L1 engages the operator fully. Placing waypoints at L2
or L3 engages the operator partly. Restarting or approving an L4 rover
engages the operator least. Let $u(t) = g_{\ell(t)}$ while the
operator interacts with a rover at level $\ell(t)$, and $u(t) = 0$
while the operator checks rovers or is idle. Engagement then evolves as
\begin{equation}
  \frac{dE}{dt} = \kappa\, u\,(1 - E) - \lambda\,(1 - u)\,E ,
  \label{eq:engagement}
\end{equation}
with gain rate $\kappa$ and decay rate $\lambda$. Checks do not raise
engagement: looking at a rover's camera feed is monitoring, not
control. When the operator switches between rovers faster than $E$
changes, Eq.~\eqref{eq:engagement} can be averaged over the switching,
and engagement settles at
\begin{equation}
  \bar{E} = \frac{\kappa\,\bar{u}}{\kappa\,\bar{u} + \lambda\,(1-\bar{u})},
  \qquad
  \bar{u} = \sum_{i=1}^{N} g_{\ell_i}\,\frac{\mathrm{IT}^{\mathrm{eff}}_i}{T_i},
  \label{eq:engagement-ss}
\end{equation}
where $T_i = \mathrm{NT}^{\mathrm{eff}}_i + \mathrm{WT}_i +
\mathrm{IT}^{\mathrm{eff}}_i$ is the length of rover $i$'s cycle, so
$\mathrm{IT}^{\mathrm{eff}}_i / T_i$ is the share of the operator's
time spent interacting with it.

An engaged operator has less to reconstruct before recovering a stuck
rover. We extend Eq.~\eqref{eq:it-rec} to
\begin{equation}
  \mathrm{IT}^{\mathrm{rec}}_\ell(a_i, E) =
  \left( \mathrm{IT}^{\mathrm{rec}}_{\ell,0} + \beta_\ell\, a_i \right)
  \left( 1 - \gamma_\ell\, E \right),
  \qquad 0 \le \gamma_\ell < 1 ,
  \label{eq:it-rec-e}
\end{equation}
where $\gamma_\ell$ is the share of a recovery at level $\ell$ that an
engaged operator can skip. We expect $\gamma_\ell$ to grow with
$\ell$, because an L4 rover's choices of where to go are the ones the
operator is least likely to have seen. With $\gamma_\ell = 0$,
Eq.~\eqref{eq:it-rec-e} reduces to Eq.~\eqref{eq:it-rec} and the
rovers are independent again. With $\gamma_\ell > 0$ they are
coupled: a rover held at a low level consumes more of the operator's
time, but it raises $\bar{u}$ and therefore $\bar{E}$, and that
shortens the recovery of every other rover in the fleet.

\subsection{Fan-Out of a Mixed Fleet}
\label{sec:fo-mixed}
Let $\boldsymbol{\ell} = (\ell_1, \dots, \ell_N)$ assign a level to
each rover in a fleet of $N$, and let $\mathcal{L}_i$ be the set of
levels feasible for rover $i$; for example, L3 is feasible only for a
rover whose waypoints lie inside the mapped region
(Section~\ref{sec:levels}). Rovers may face different tasks, so we
write the parameters of rover $i$ at its assigned level with subscript
$i$ (e.g., $\mathrm{NT}_i$, $h_i$, $\gamma_i$); they may depend on the
rover's task as well as its level.

Eq.~\eqref{eq:fo-aware} has no meaning in a mixed fleet, because there
is no single level to evaluate it at. Instead we count operator time.
In each cycle, rover $i$ takes $\mathrm{IT}^{\mathrm{eff}}_i$ of
interaction and $\tau_i$ of trust tax, so its share of the operator's
time is $\rho_i = (\mathrm{IT}^{\mathrm{eff}}_i + \tau_i)/T_i$, and the
operator's utilization is
\begin{equation}
  U(\boldsymbol{\ell}) = \sum_{i=1}^{N} \rho_i
  = \sum_{i=1}^{N} \frac{\mathrm{IT}^{\mathrm{eff}}_i + \tau_i}
                         {\mathrm{NT}^{\mathrm{eff}}_i + \mathrm{WT}_i + \mathrm{IT}^{\mathrm{eff}}_i} .
  \label{eq:utilization}
\end{equation}
For a uniform fleet with no trust tax and no wait, setting $U = 1$
gives $N = \mathrm{NT}/\mathrm{IT} + 1$, which is
Eq.~\eqref{eq:fo-classic} with $\mathrm{WT} = 0$. An operator cannot
sustain full utilization over a shift; performance degrades once
utilization exceeds roughly 70\%~\cite{cummings2007developing}. We
therefore define the fan-out of a mixed fleet as the largest fleet
that some feasible assignment keeps within a sustainable utilization
$U_{\max} < 1$:
\begin{equation}
  \mathrm{FO}^{*} = \max \Big\{ N : \min_{\boldsymbol{\ell} \in
  \mathcal{L}_1 \times \dots \times \mathcal{L}_N}
  U(\boldsymbol{\ell}) \le U_{\max} \Big\} .
  \label{eq:fo-mixed}
\end{equation}

$U(\boldsymbol{\ell})$ is not a sum of independent terms. Each
$\rho_i$ depends on $\bar{E}$ through rover $i$'s recovery time
(Eq.~\eqref{eq:it-rec-e}); $\bar{E}$ depends on $\bar{u}$, which
depends on every rover's interaction time
(Eq.~\eqref{eq:engagement-ss}); and the queueing delay $W_q$ inside
each $\mathrm{WT}_i$ grows with $U$ itself. We evaluate $U$ for a
given assignment by fixed-point iteration over
Eqs.~\eqref{eq:wt}, \eqref{eq:it-eff}, \eqref{eq:engagement-ss},
and~\eqref{eq:it-rec-e}, starting from $\bar{E} = 0$.
\todo{Check convergence of the iteration for the fitted parameters.}

\emph{When mixing helps.} Lowering rover $k$'s level changes
utilization by approximately
\begin{equation}
  \Delta U \approx \Delta \rho_k
  + \Delta \bar{E} \sum_{j \ne k} \frac{\partial \rho_j}{\partial \bar{E}} .
  \label{eq:mix-condition}
\end{equation}
The first term is usually positive, since a rover at a lower level
needs the operator more often. The second is negative, since
$\Delta \bar{E} > 0$ and engagement shortens recovery
($\partial \rho_j / \partial \bar{E} \le 0$). Holding rover $k$ at a
lower level pays off when the engagement it produces, summed over the
other $N-1$ rovers, saves more operator time than rover $k$ costs.
Each term in the sum scales with $h_j \gamma_j$, the rate at which
rover $j$ fails times the share of its recovery that engagement
removes. The model therefore predicts that mixing levels helps most in
large fleets of highly autonomous rovers that fail often in ways the
operator cannot see coming, and that it does not help in small fleets
or with reliable autonomy, where a uniform high level is best.

\emph{Allocation and engagement gains.} Let $\boldsymbol{\ell}^{U}$ be
the best uniform assignment, $\boldsymbol{\ell}^{0}$ the best
assignment under the model with $\gamma_i = 0$ for all rovers, and
$\boldsymbol{\ell}^{*}$ the best assignment under the full model. With
all three evaluated under the full model, the saving in utilization
from mixing splits as
\begin{equation}
  U(\boldsymbol{\ell}^{U}) - U(\boldsymbol{\ell}^{*}) =
  \underbrace{U(\boldsymbol{\ell}^{U}) - U(\boldsymbol{\ell}^{0})}_{\text{allocation}}
  + \underbrace{U(\boldsymbol{\ell}^{0}) - U(\boldsymbol{\ell}^{*})}_{\text{engagement}} .
  \label{eq:decomposition}
\end{equation}
The allocation gain comes from matching each rover's level to its
task while treating rovers as independent. The engagement gain comes
from choosing levels with the coupling through $\bar{E}$ in view, and
it is non-negative by construction because $\boldsymbol{\ell}^{*}$ is
optimal under the full model. If $\gamma_i = 0$ for every rover, the
engagement gain is zero, so an empirical estimate of $\gamma_\ell$ is
a direct test of the engagement claim.

\subsection{Autonomy Allocation}
\label{sec:allocation}
During the task, the optimizer recommends a level for each rover.
After each observed outcome, it updates each $\theta_i$ with
Eq.~\eqref{eq:trust} and the engagement estimate $\hat{E}$ with
Eq.~\eqref{eq:engagement} from the logged interactions, and solves
\begin{equation}
  \boldsymbol{\ell}^{+} = \arg\min_{\boldsymbol{\ell} \in
  \mathcal{L}_1 \times \dots \times \mathcal{L}_N}
  \; U(\boldsymbol{\ell};\, \boldsymbol{\theta}, \hat{E})
  + \frac{1}{H} \sum_{i=1}^{N} s\big(\ell_i^{\mathrm{cur}} \to \ell_i\big),
  \label{eq:allocation}
\end{equation}
where $s(\ell \to \ell')$ is the operator time needed to move a rover
from its current level $\ell$ to $\ell'$ (for example, taking over a
rover the operator has not been watching), amortized over a planning
horizon $H$. To avoid oscillating between assignments, the optimizer
recommends a change only when it lowers $U$ by more than a margin
$\mu$. The operator may accept or decline each recommendation. We log
every recommendation and response; a declined recommendation is a
further behavioral signal of trust. With five levels the search space
has at most $5^N$ assignments, which we search exhaustively for the
fleet sizes in our study \todo{state fleet sizes}. For larger fleets,
because the assignment enters $U$ only through per-rover terms and the
scalar $\bar{u}$, coordinate descent over rovers is a natural
approximation. \todo{State how the model parameters are estimated,
e.g., from pilot sessions.}



\section*{Acknowledgment}

\bibliographystyle{IEEEtran}
\bibliography{references}

\end{document}
